\section{Conclusion}

We began by noting that machine learning benchmarks drive billions in investment without systematic saturation detection mechanisms. Our temporal precedence analysis reveals that saturation signals appear an average of 4 years before community migration---far longer than previously assumed, providing an actionable early warning window for proactive benchmark rotation.

Through proof-of-concept experiments on major benchmarks (ImageNet, GLUE, SQuAD), we demonstrated three key contributions. First, expert consensus on saturation timing exists with >70\% agreement (76-93\% within $\pm$1 year of modal dates), validating saturation as community-observable phenomenon. Second, score convergence detection via rolling window statistics achieves 100\% detection rate with statistical significance (Levene's $p<0.05$), demonstrating algorithmic detectability without complex ML models. Third, 100\% of tested saturations preceded paradigm shift adoption by 32-78 months (mean 48 months), distinguishing internal benchmark exhaustion from external disruption.

These findings provide validated mechanisms that, pending real-world deployment, could enable a shift from reactive organic migration to systematic benchmark lifecycle management. Treating benchmarks as consumables with time-boxed validity periods (not permanent monuments) provides the foundation for saturation detection infrastructure. Conference organizers can integrate saturation dashboards into submission systems, enabling benchmark selection rationale disclosure and voluntary adoption.

\subsection{Future Work}

\textbf{Immediate Extensions (High Priority)}:
\begin{itemize}
\item \textbf{Real PWC Data Validation (FW-7)}: Re-run all experiments with actual Papers With Code leaderboard data when API access restored. Unlocks real-world applicability claims and determines whether per-benchmark calibration (L4) is genuinely required or synthetic artifact.
\item \textbf{Dual-Metric Integration (FW-1)}: Implement combined detector requiring BOTH score convergence AND velocity decay. Compare precision/recall against single-metric baselines on expanded benchmark set ($n=10+$).
\item \textbf{Real Expert Survey Collection (FW-4)}: Execute full expert survey with IRB approval, targeting $n=100$-150 responses via NeurIPS/ICML mailing lists. Validate assumption A1 (expert consensus exists) with actual ML researcher opinions.
\end{itemize}

\textbf{Sample Size Expansion (Medium Priority)}:
\begin{itemize}
\item \textbf{Expanded Temporal Precedence Sample (FW-8)}: Grow from $n=3$ to $n=15+$ benchmark-shift pairs for statistical significance (current $p=0.125 \rightarrow$ target $p<0.05$). Additional pairs: CIFAR-10$\rightarrow$ResNet, SuperGLUE$\rightarrow$T5, WikiText$\rightarrow$GPT-2, MS COCO$\rightarrow$Mask R-CNN.
\end{itemize}

\textbf{Citation Mechanism Refinement (Medium Priority)}:
\begin{itemize}
\item \textbf{Citation Correlation Window Alignment (FW-2)}: Test aligned 9-month windows (both detector and ground truth) and threshold variants ($3\sigma$, $4\sigma$ spike detection) to improve precision from 0.50 toward 0.80 target.
\end{itemize}

\textbf{Longer-Term Vision}:
\begin{itemize}
\item \textbf{Benchmark Governance Protocols}: Develop community coordination mechanisms for systematic benchmark lifecycle management.
\item \textbf{Conference Policy Pilots (FW-6)}: Integrate saturation dashboard with NeurIPS Benchmark Track or ICLR to test voluntary adoption hypothesis. Track adoption rate (\% papers referencing saturation status) and author survey feedback.
\end{itemize}

The shortest path from proof-of-concept to production deployment is real PWC data validation (FW-7), which determines whether mechanisms validated on synthetic data generalize to actual leaderboard submissions. This single validation unlocks applicability claims and enables infrastructure piloting with conference organizers.

Our work demonstrates that benchmark saturation is not an inevitable background process to be endured, but a detectable phenomenon with measurable signals and actionable lead times. Systematic saturation detection transforms benchmark lifecycle from unmanaged monument to governed consumable, enabling the ML research community to invest resources where genuine progress remains possible.
